\begin{abstract}
Continual Learning in Visual Question Answering (VQACL) requires models to acquire new visual-linguistic skills (plasticity) while preserving previously learned knowledge (stability). The inherent multimodality of VQACL exacerbates this challenge, as models must balance stability across visual and textual domains while adapting to novel objects and reasoning tasks. Existing methods, primarily designed for unimodal settings, often fall short in addressing this dual requirement. In this work, we present \textit{QUestion-only replay with Attention Distillation} (\qstmethodshort{}), a novel approach for VQACL that leverages only past task questions for regularization. By eliminating the need to store visual data, \qstmethodshort{} not only reduces memory overhead, but also alleviates privacy concerns. Our method introduces a \textit{Question-only Replay} mechanism that selectively reuses prior task questions to counteract overfitting to the answer space of the current task, addressing the problem \textit{ out of answer set}. Complementing this, we propose \textit{Attention Consistency Distillation} to enforce both intra-modal and inter-modal attention consistency across tasks, preserving essential visual-linguistic associations. Extensive experiments on VQAv2 and NExT-QA demonstrate that \qstmethodshort{} significantly outperforms state-of-the-art methods, achieving robust performance in continual VQA. Code is available at: \href{https://github.com/IemProg/QUAD}{https://github.com/IemProg/QUAD}. \footnote{Work done during an internship at Computer Vision Center (CVC),
Universitat Autonoma de Barcelona.}
\end{abstract}